\documentclass[12pt]{article}
\usepackage[a4paper,margin=1in]{geometry}
\usepackage[utf8]{inputenc}
\usepackage[T1]{fontenc}
\usepackage{amsmath,amssymb}
\usepackage{textcomp}
\usepackage{parskip}
\begin{document}
\begin{center}
{\Large \textbf{2025 STEP II}}
\end{center}

\subsection*{Page 1}
\begin{flushleft}
\bigskip\textbf{Section A: Pure Mathematics}\par
\bigskip\textbf{Question 1.} The function Min is defined as\par
a if a  b\par
\[Min(a, b) =\]
\[b if a > b .\]
\textbf{(i)} Sketch the graph y = Min(x2, 2x).\par
\textbf{(ii)} Solve the equation 2 Min(x2, 2x) = 5x 3.\par
-\par
\textbf{(iii)} Solve the equation Min(x2, 2x) + Min(x3, 4x) = mx in the cases m = 2 and m = 6.\par
\textbf{(iv)} Show that (1, 3) is a local maximum point on the curve y = 2 Min(x2, x3) 5x and\par
- -\par
find the other three local maxima and minima on this curve.\par
Sketch the curve.\par
\end{flushleft}
\subsection*{Page 2}
\begin{flushleft}
2 (i) (a) Show that if the complex number z satisfies the equation\par
\[z^2 + z + b = a,\]
| |\par
where a and b are real numbers, then z must be either purely real or purely\par
imaginary.\par
(b) Show that the equation\par
\[z^2 + z + 5 = 7\]
2 2\par
has no purely imaginary roots.\par
(c) Show that the equation\par
\[z^2 + z + 7 = 5\]
2 2\par
has no purely real roots.\par
\[(d) Show that, when 1 < b < 3 , the equation\]
2 4\par
\[z^2 + z + b = 1\]
| | 2\par
will have at least one purely imaginary root and at least one purely real root.\par
\textbf{(ii)} Solve the equation\par
\[z^3 + z + 2 2 = 4.\]
| |\par
ln x\par
\bigskip\textbf{Question 3.} (i) Sketch a graph of y = for x > 0.\par
x\par
\textbf{(ii)} Use your graph to show the following.\par
\[(a) 3\pi > \pi3\]
\[(b) 9 \sqrt\{5\} > \sqrt\{5\} 9 4\]
( )\par
\textbf{(iii)} Given that 1 < x < 2, decide, with justification, which is the larger of xx+2 or (x + 2)x.\par
\textbf{(iv)} Show that the inequalities 9\sqrt2 > \sqrt2 9 and 32\sqrt2 > 2\sqrt2 3 are equivalent. Given that\par
\[e^2 < 8, decide, with justification, which is the larger of 9\sqrt\{2\} and \sqrt\{2\} 9 .\]
( )\par
\textbf{(v)} Decide, with justification, which is the larger of 8 \sqrt3 3 and \sqrt3 3 8 .\par
\end{flushleft}
\subsection*{Page 3}
\begin{flushleft}
\bigskip\textbf{Question 4.} Let x denote the largest integer that satisfies x  x.\par
(  (\par
\[For example, if x = 4.2, then x = 5.\]
-  ( -\par
\textbf{(i)} Show that, if n is an integer, then x + n = x + n.\par
(  (\par
\textbf{(ii)} Let n be a positive integer and define function f by\par
n\par
1 2 n 1\par
\[f (x) = x + x + + x + + . . . + x + - nx\]
n\par
( n n n -  (\par
(  (  (\par
\[(a) Show that f x + = f (x).\]
n n\par
n\par
) )\par
\[(b) Evaluate f n (t) for 0  t < .\]
n\par
(c) Hence show that f (x) 0.\par
n\par
\equiv\par
x x + 1\par
\textbf{(iii)} (a) Show that + = x .\par
2 2  (\par
(  (\par
(b) Hence, or otherwise, simplify\par
x + 1 x + 2 x + 2k\par
+ + . . . + + . . . .\par
2 22 2k+1\par
(  (  (\par
\end{flushleft}
\subsection*{Page 4}
\begin{flushleft}
\bigskip\textbf{Question 5.} You need not consider the convergence of the improper integrals in this question.\par
\textbf{(i)} Use the substitution x = u 1 to show that\par
-\par
\[\sqrt\{x\} 1\]
\infty\par
\[- dx = 0 .\]
\[x(x^3 + 1)\]
(\par
\textbf{(ii)} Use the substitution x = u 2 to show that\par
-\par
1 1\par
\infty \infty\par
\[dx = 2 dx .\]
\[\sqrt\{x^3\} + 1 \sqrt\{x^6\} + 1\]
0 0\par
\textbf{(iii)} Find, in terms of p and s, a value of r for which\par
xr 1\par
\infty\par
\[- dx = 0,\]
xs(xp + 1)\par
given that p and s are fixed values (for which the required integrals converge.\par
\textbf{(iv)} Show that, for any positive value of k, it is possible to find values of p and q for which\par
1 1\par
\infty \infty\par
\[dx = k dx .\]
\[\sqrt\{xp\} + 1 \sqrt\{xq\} + 1\]
0 0\par
\end{flushleft}
\subsection*{Page 5}
\begin{flushleft}
\bigskip\textbf{Question 6.} (i) The circle x2 + (y a)2 = r2 touches the parabola 2ky = x2, where k > 0, tangentially\par
-\par
\[at two points. Show that r^2 = k(2a k).\]
-\par
\[Show further that if r^2 = k(2a k) and a > k > 0, then the circle x^2 + (y a)2 = r^2\]
- -\par
\[touches the parabola 2ky = x^2 tangentially at two points.\]
\textbf{(ii)} The lines y = c x are tangents to the circle x2 + (y a)2 = r2. Find r2, and the\par
± -\par
coordinates of the points of contact, in terms of a and c.\par
\textbf{(iii)} C and C are circles with equations x2 + (y a )2 = r2 and x2 + (y a )2 = r2\par
1 2 - 1 1 - 2 2\par
\[respectively, where a = a and r = r .\]
1 2 1 2\par
) )\par
\[Each circle touches the parabola 2ky = x^2 tangentially at two points and the lines\]
\[y = c x are tangents to both circles.\]
±\par
\[(a) Show that a + a = 2c + 4k and that a^2 + a^2 = 2c2 + 16kc + 12k2.\]
1 2 1 2\par
\[(b) The circle x^2 + (y d)2 = p^2 passes through the four points of tangency of the\]
-\par
\[lines y = c x to the two circles, C and C . Find d and p^2 in terms of k and c.\]
1 2\par
±\par
\[(c) Show that the circle x^2 + (y d)2 = p^2 also touches the parabola 2ky = x^2\]
-\par
tangentially at two points.\par
\bigskip\textbf{Question 7.} The differential equation\par
\[d2x dx\]
\[= 2x\]
\[dt2 dt\]
describes the motion of a particle with position x(t) at time t. At t = 0, x = a, where a > 0.\par
dx\par
\textbf{(i)} Solve the differential equation in the case where = a2 when t = 0.\par
dt\par
What happens to the particle as t increases from 0?\par
dx\par
\textbf{(ii)} Solve the differential equation in the case where = a2 + p2 when t = 0, where p > 0.\par
dt\par
What happens to the particle as t increases from 0?\par
dx\par
\textbf{(iii)} Solve the differential equation in the case where = a2 q2 when t = 0, where q > 0.\par
dt -\par
What happens to the particle as t increases from 0? Give conditions on a and q for\par
the different cases which arise.\par
\end{flushleft}
\subsection*{Page 6}
\begin{flushleft}
8 If we split a set S of integers into two subsets A and B whose intersection is empty and\par
whose union is the whole of S, and such that\par
the sum of the elements of A is equal to the sum of the elements of B\par
•\par
and the sum of the squares of the elements of A is equal to the sum of the squares of\par
•\par
the elements of B,\par
then we say that we have found a balanced partition of S into two subsets.\par
\textbf{(i)} Find a balanced partition of the set 1, 2, 3, 4, 5, 6, 7, 8 into two subsets A and B, each\par
\{ \}\par
of size 4.\par
\textbf{(ii)} Given that a , a , . . . , a and b , b , . . . , b are sequences with\par
1 2 m 1 2 m\par
m m m m\par
\[a = b and a^2 = b^2,\]
k k k k\par
\[k=1 k=1 k=1 k=1\]
show that\par
m m m m\par
\[a^3 + (c + b )3 = b^3 + (c + a )3\]
k k k k\par
\[k=1 k=1 k=1 k=1\]
for any real number c.\par
\textbf{(iii)} Find, with justification, a balanced partition of the set 1, 2, 3, . . . , 16 into two subsets\par
\{ \}\par
A and B, each of size 8, which also has the property that\par
the sum of the cubes of the elements of A is equal to the sum of the cubes of the\par
•\par
elements of B.\par
\textbf{(iv)} You are given that the sets A = 1, 3, 4, 5, 9, 11 and B = 2, 6, 7, 8, 10 form a balanced\par
\{ \} \{ \}\par
partition of the set 1, 2, 3, . . . , 11 .\par
\{ \}\par
\[Let S = n^2, (n + 1)2, (n + 2)2, . . . , (n + 11)2 , where n is any positive integer. Find,\]
\{ \}\par
with justification, two subsets C and D of S whose intersection is empty and whose\par
union is the whole of S, and such that\par
the sum of the elements of C is equal to the sum of the elements of D.\par
•\par
\end{flushleft}
\subsection*{Page 7}
\begin{flushleft}
\bigskip\textbf{Section B: Mechanics}\par
\bigskip\textbf{Question 9.} Points A and B are at the same height and a distance \sqrt2 r apart. Two small, spherical\par
particles of equal mass, P and Q, are suspended from A and B, respectively, by light inex-\par
tensible strings of length r. Each particle individually may move freely around and inside a\par
circle centred at the point of suspension.\par
The particles are projected simultaneously from points which are a distance r vertically below\par
their points of suspension, directly towards each other and each with speed u. When the\par
particles collide, the coefficient of restitution in the collision is e.\par
\textbf{(i)} Show that, immediately after the collision, the horizontal component of each particle's\par
\[velocity has magnitude 1 ev\sqrt\{2\}, where v^2 = u^2 gr 2 \sqrt\{2\} and write down the\]
2 - -\par
vertical component in terms of v.\par
(\par
\textbf{(ii)} Show that the strings will become taut again at a time t after the collision, where t is\par
a non-zero root of the equation\par
\[(r evt)2 + r + vt 1 \sqrt\{2gt2\} = 2r2.\]
- - - 2\par
) )\par
\textbf{(iii)} Show that, in terms of the dimensionless variables\par
\[vt \sqrt\{2v2\}\]
\[z = and c =\]
r rg\par
this equation becomes\par
z 3 z 2 2 z 2\par
\[2 + + 1 + e^2 (1 + e) = 0.\]
c - c c c - c\par
( )\par
) ) ) ) ) )\par
\textbf{(iv)} Show that, if this equation has three equal non-zero roots, e = 1 and v2 = 9 \sqrt2rg.\par
3 2\par
Explain briefly why, in this case, no energy is lost when the string becomes taut.\par
\textbf{(v)} In the case described in (iv), the particles have speed U when they again reach the\par
points of their motion vertically below their points of suspension. Find U 2 in terms of\par
r and g.\par
\end{flushleft}
\subsection*{Page 8}
\begin{flushleft}
\bigskip\textbf{Question 10.} The lower end of a rigid uniform rod of mass m and length a rests at point M on rough\par
horizontal ground. Each of two elastic strings, of natural length ) and modulus of elasticity\par
λ, is attached at one end to the top of the rod. Their lower ends are attached to points A\par
and B on the ground, which are a distance 2a apart. M is the midpoint of AB.\par
P is the point at the top of the rod and lies in the vertical plane through AM B.\par
\[Suppose that the rod is in equilibrium with angle P M B = 2θ, where θ < 45 and ) is such\]
\textasciicircum{}\circ\par
that both strings are in tension.\par
\textbf{(i)} Show that angle AP B is a right angle.\par
Show that that the force exerted on the rod by the elastic strings can be written as\par
the sum of\par
2aλ\par
a force of magnitude parallel to the rod\par
• )\par
\[and a force of magnitude \sqrt\{2\} λ acting along the bisector of angle AP B.\]
•\par
2λ\par
\textbf{(ii)} By taking moments about point M , or otherwise, show that cos θ + sin θ = .\par
mg\par
\[Deduce that it is necessary that 1 mg < λ < 1 \sqrt\{2\} mg.\]
2 2\par
\textbf{(iii)} N and F are the magnitudes of the normal and frictional forces, respectively, exerted\par
on the rod by the ground at M .\par
Show, by taking moments about an appropriate point, or otherwise, that\par
\[N F tan 2θ = 1 mg.\]
- 2\par
\end{flushleft}
\subsection*{Page 9}
\begin{flushleft}
\bigskip\textbf{Section C: Probability and Statistics}\par
\bigskip\textbf{Question 11.} (i) By considering the sum of a geometric series, or otherwise, show that\par
\infty 1\par
\[rxr 1 = for x < 1.\]
-\par
(1 x)2 | |\par
\[r=1 -\]
\textbf{(ii)} Ali plays a game with a fair 2k-sided die. He rolls the die until the first 2k appears.\par
Ali wins if all the numbers he rolls are even.\par
(a) Find the probability that Ali wins the game.\par
If Ali wins the game, he earns \pounds 1 for each roll, including the final one. If he loses, he\par
earns nothing.\par
(b) Find Ali's expected earnings from playing the game.\par
\textbf{(iii)} Find a simplified expression for\par
n n\par
\[1 + 2 x + 3 x^2 + . . . + (n + 1)xn,\]
1 2\par
( ) ( )\par
where n is a positive integer.\par
\textbf{(iv)} Zen plays a different game with a fair 2k-sided die. She rolls the die until the first 2k\par
appears, and wins if the numbers rolled are strictly increasing in size. For example, if\par
k = 3, she wins if she rolls 2, 6 or 1, 4, 5, 6, but not if she rolls 1, 4, 2, 6 or 1, 3, 3, 6.\par
If Zen wins the game, she earns \pounds 1 for each roll, including the final one. If she loses,\par
she earns nothing.\par
Find Zen's expected earnings from playing the game.\par
\textbf{(v)} Using the approximation\par
n\par
1 + e for large n,\par
n \approx\par
( )\par
show that, when k is large, Zen's expected earnings are a little over 35\% more than\par
Ali's expected earnings.\par
\end{flushleft}
\subsection*{Page 10}
\begin{flushleft}
\bigskip\textbf{Question 12.} Let X be a Poisson random variable with mean λ and let p = P(X = r), for r = 0, 1, 2, . . . .\par
r\par
Neither λ nor λ + 1 + λ + 1 is an integer.\par
2 4\par
\textbf{(i)} Show, by considering the sequence d p p for r = 1, 2, . . . , that there is a\par
r r r 1\par
\equiv - -\par
\[unique integer m such that P(X = r)  P(X = m) for all r = 0, 1, 2, . . . , and that\]
\[λ 1 < m < λ.\]
-\par
\textbf{(ii)} Show that the minimum value of d occurs at r = k, where k is such that\par
r\par
\[k < λ + 1 + λ + 1 < k + 1.\]
2 4\par
\textbf{(iii)} Show that the condition for the maximum value of d to occur at r = 1 is\par
r\par
\[1 < λ < 2 + \sqrt\{2\}.\]
\textbf{(iv)} In the case λ = 3.36, sketch a graph of p against r for r = 0, 1, 2, . . . , 6, 7.\par
r\par
\end{flushleft}
\end{document}
